\Needspace{18\baselineskip}
\section{Worked problems}

\Needspace{12\baselineskip}
\subsection{Binary trees}
\textbf{Problem.} \emph{Tree arithmetic.} Compute node count and maximum accepted draft length for
full binary trees of depths 1 through 5. Explain why node count is not acceptance length.

\textbf{Solution.} A full binary tree of depth $d$ contains
$N=2^{d+1}-2$ draft nodes when the committed root is excluded, but one accepted path
contains at most $d$ draft tokens. For $d=1,\ldots,5$, $N=2,6,14,30,62$.

\Needspace{12\baselineskip}
\subsection{Parent array}
\textbf{Problem.} \emph{Mask construction.} For parent array $[-1,0,0,1,2]$, write the five-by-five
draft visibility matrix and logical depth of every node.

\textbf{Solution.} For $[-1,0,0,1,2]$, depths are $(1,2,2,3,3)$ if node 0 is
the first draft after the committed prefix. Visibility sets are
$A:\{A\}$, $B:\{A,B\}$, $C:\{A,C\}$, $D:\{A,B,D\}$, $E:\{A,C,E\}$.
The rows for C and E expose why a triangular mask is wrong.

\Needspace{12\baselineskip}
\subsection{Smallest distinguishing tree}
\textbf{Problem.} \emph{Differential test.} Explain the smallest tree that distinguishes ancestry
masking from ordinary triangular masking.

\textbf{Solution.} One parent with two children suffices.
Flatten $[A,B,C]$ with both B and C children of A. A triangular mask lets C see B;
ancestry masking does not.

\Needspace{12\baselineskip}
\subsection{Budget priority}
\textbf{Problem.} \emph{Budgeting.} Given frontier scores/costs $(.8,2),(.5,1),(.9,3),(.3,1)$,
rank them by $\rho=p/c$. Explain one reason that greedy ranking may still be suboptimal.

\textbf{Solution.} Scores $p/c$ are $.4,.5,.3,.3$, so the second item ranks
first, then first, then the tied remainder. Greedy density can be suboptimal because
expanding one node reveals valuable descendants and because verification cost is not
strictly additive across GPU shapes.
